%% A Synthetic Journal-Bearing Benchmark for Stress-Testing
%% Physics-Informed Fault Diagnosis  (artifact: JBFD-11)
%%
%% Manuscript style: NeurIPS Datasets & Benchmarks / Evaluations.
%% Build:  latexmk -pdf main.tex     (offline; self-contained neurips_2023.sty)
%% arXiv categories: cs.LG (primary) + eess.SP (secondary); cross-list stat.ML.
%%
%% BINDING CLAIM BOUNDARIES (paper/OUTLINE.md S5; results/FINDINGS.md S0):
%%   - NO "physics improves ..." wording of any kind (not even "a spectral
%%     regularizer helped"). Complete negative on the TESTED methods only.
%%   - Record-level (528) + seed-level estimands ONLY. The within-seed /
%%     representative-seed McNemar (14--0, p=1.2e-4) is reported ONLY as the
%%     cautionary confounded statistic, never as evidence of an effect.
%%   - Synthetic-only, stated throughout. "will archive" until the Zenodo DOI.
%%   - Every number traces to an artifact path under results/.

\documentclass{article}

% [preprint] = authors shown, no line numbers (arXiv). Switch to no-option for
% an anonymous D&B submission, or [final] for camera-ready.
\usepackage[preprint]{neurips_2023}

\usepackage[utf8]{inputenc}
\usepackage[T1]{fontenc}
\usepackage{hyperref}
\usepackage{url}
\usepackage{booktabs}
\usepackage{amsfonts}
\usepackage{amsmath}
\usepackage{nicefrac}
\usepackage{microtype}
\usepackage{graphicx}
\usepackage{xcolor}

\graphicspath{{figures/}}

% Include a figure if its file exists, else a labelled placeholder box, so the
% skeleton always compiles while figures are still being generated.
\newcommand{\figfromfile}[1]{%
  \IfFileExists{figures/#1}%
    {\includegraphics[width=\linewidth]{#1}}%
    {\fbox{\parbox[c][3.2cm][c]{0.98\linewidth}{\centering\itshape
       figure placeholder --- \texttt{figures/#1} not yet generated}}}%
}

\title{A Synthetic Journal-Bearing Benchmark for Stress-Testing\\
       Physics-Informed Fault Diagnosis}

% Affiliation provisional (owner: PIEAS; currently with Cowlar Design Studio);
% to be finalized before the arXiv posting. Anonymize only for a double-blind
% D&B submission.
\author{%
  Syed Abbas Ahmad\thanks{Work conducted with Cowlar Design Studio.
  Affiliation provisional; finalized before arXiv posting.}\\
  Pakistan Institute of Engineering and Applied Sciences (PIEAS)\\
  \texttt{bsso.cowlar@gmail.com}%
}

\begin{document}
\maketitle

\begin{abstract}
Physics-informed neural networks are increasingly proposed for bearing fault
diagnosis, but controlled, leakage-free evidence that physics priors actually
\emph{help} is scarce --- and the journal/hydrodynamic-bearing regime, unlike
the rolling-element datasets that dominate the literature (CWRU, Paderborn), has
almost no public benchmarks. We release \textbf{JBFD-11}, a physics-grounded
\textbf{synthetic} journal-bearing vibration dataset (11 fault classes,
severity-stratified, record-level leakage-checked, with SNR-20/10/5 variants)
from a generator whose spectral signatures are locked by a 34-test CI battery,
together with a \textbf{frozen benchmark of 11 classical, deep, and
physics-informed models, analyzed at the record level}. We use this controlled
testbed to stress-test the physics-informed methods we implement --- a
band-energy spectral-consistency loss judged against a frozen healthy-class
reference, with per-sample operating conditions --- across the regimes where
such priors are theorized to help: noise robustness, data-efficiency, severity
out-of-distribution, interpretability, and calibration. On every axis,
\textbf{the tested physics-informed methods show no statistically supported,
seed-robust advantage} over strong data-driven baselines. We further document
how readily an apparent benefit arises and then disappears under scrutiny: a
promising $n{=}3$ analysis --- a highly significant \emph{representative-seed}
McNemar result ($14$--$0$, $p=1.2\times10^{-4}$) --- \textbf{dissolves under a
pre-registered $n{=}12$ \emph{seed-level} test} (Wilcoxon $p=0.79$); a
pre-registered control that swaps the physics targets for \textbf{random
non-fault frequency bands} has the lowest mean degradation numerically yet still
fails the pre-registered robustness criterion. Our contribution is a reusable
physics-grounded benchmark, a rigorous complete negative on the tested methods,
and a concrete methodological caution about seed counts and estimands in
physics-informed evaluation. We release the dataset, generator, and analysis
code, and \textbf{will} archive the trained checkpoints and full reproducibility
package on a public repository; every headline number was independently
recomputed from the raw checkpoints by separate audit scripts.
\emph{(Synthetic-only; no real-rig validation.)}
\end{abstract}

% =====================================================================
\section{Introduction}
\label{sec:intro}
\textit{[TODO C-intro] The gap (journal-bearing data scarcity; physics-informed
diagnosis proposed but under-evaluated~\citep{raissi2019pinn}); what a controlled
synthetic testbed buys; contributions; the headline negative and the
$n{=}3\!\rightarrow\!n{=}12$ hook (Fig.~\ref{fig:dissolution}). State
synthetic-only here. Rolling-element benchmarks dominate the
literature~\citep{smith2015cwru,lessmeier2016paderborn}.}

% =====================================================================
\section{Related work}
\label{sec:related}
\textit{[TODO] Vibration diagnosis (signal processing $\rightarrow$ deep
learning); PINNs / physics-informed diagnosis~\citep{raissi2019pinn};
synthetic vibration data; evaluation rigor, negative results, and
reproducibility in ML.}

% =====================================================================
\section{Dataset \& generator [C1]}
\label{sec:dataset}

\paragraph{Motivation and scope.}
JBFD-11 is a physics-grounded \emph{synthetic} vibration dataset for
hydrodynamic journal-bearing fault diagnosis. Unlike rolling-element bearings,
for which public benchmarks are abundant (CWRU~\citep{smith2015cwru},
Paderborn~\citep{lessmeier2016paderborn}), journal/hydrodynamic bearings have
almost no openly available labelled fault data, yet they fail through distinct
mechanisms with distinct vibration signatures. We therefore generate a
controlled testbed in which the ground-truth fault, its severity, and the
operating conditions are known exactly, so that an evaluation can isolate
\emph{whether a physics prior helps} rather than confounding it with
uncontrolled field variability. The dataset is synthetic throughout; we make no
claim of real-rig fidelity (\S\ref{sec:limitations}).

\paragraph{Signal model.}
Each record is generated as $x(t) = n_{\mathrm{base}}(t) + \sum_\ell
\mathrm{noise}_\ell(t) + s_{\mathrm{fault}}(t)$, sampled at $f_s=20{,}480$\,Hz
for $5$\,s ($102{,}400$ samples) at a nominal shaft speed $\Omega_0=60$\,Hz
($3{,}600$\,RPM). A machine floor $n_{\mathrm{base}}$ and a seven-layer noise
model (measurement, mains EMI, pink $1/f$, thermal drift, ADC quantization,
calibration drift, sporadic impulses, and an occasional aliasing artifact) are
added to a per-class fault signature $s_{\mathrm{fault}}$ scaled by severity and
a transient envelope. Every equation is normative in \texttt{docs/PHYSICS.md}
and locked by a 34-test spectral CI battery
(\texttt{tests/test\_physics\_signatures.py}). The amplitude coefficients are
engineering-plausible values rather than Reynolds-equation solutions: the
\emph{frequency structure}, not absolute g-levels, carries the diagnostic
content.

\paragraph{Classes and signatures.}
The taxonomy comprises eight single faults --- healthy (\texttt{sain}, noise
floor only), misalignment (2X/3X harmonics), imbalance (1X growing with
speed$^2$), clearance/looseness (sub-synchronous $0.43$--$0.48{\times}$ plus
1X/2X), lubrication deficiency (2--5\,Hz stick-slip, amplitude $\propto 1/S$ in
the Sommerfeld number $S$), cavitation (1.4--2.6\,kHz bursts), wear (broadband
floor plus asperity harmonics), and oil whirl ($0.42$--$0.48{\times}$, amplitude
$\propto 1/\sqrt{S}$) --- and three mixed faults that superpose constituents at
reduced amplitude (Fig.~\ref{fig:dataset}). Each signature follows established
vibration-analysis practice, cited per class in \texttt{docs/PHYSICS.md}\,\S4.
Cavitation is only weakly expressed spectrally; we flag this and qualify any
cavitation-specific reading.

\begin{figure}[t]
  \centering
  \figfromfile{F3_dataset_overview.pdf}
  \caption{JBFD-11 signature map: the 11 classes against their characteristic
    vibration signatures at the nominal shaft speed $\Omega_0=60$\,Hz, with the
    tonal/shaft-order and broadband/HF families distinguished. Frequencies shown
    are nominal centres for visualization; the normative equations are in
    \texttt{docs/PHYSICS.md}\,\S4.}
  \label{fig:dataset}
\end{figure}

\paragraph{Severity, operating conditions, and physics couplings.}
Faults are stratified into four non-overlapping severity levels (incipient
$0.20$--$0.45$, mild $0.45$--$0.70$, moderate $0.70$--$0.90$, severe
$0.90$--$1.00$) with exactly 80 records per class per level; 30\% of records
carry a linearly progressing severity envelope. Operating conditions are sampled
per record and stored in metadata: shaft speed $\Omega\in[54,66]$\,Hz
($\pm 10\%$), load 30--100\% of rated, and temperature 40--80$^\circ$C, from
which an oil viscosity $\mu(T)=\mu_{\mathrm{ref}}e^{-0.03(T-60)}$ and a
Sommerfeld number
$S=S_{\mathrm{base}}\,(\mu/\mu_{\mathrm{ref}})(\Omega/\Omega_0)/\mathrm{load}$
(clipped to $[0.05,0.50]$) drive the lubrication and oil-whirl couplings. The
recorded per-record conditions provide the per-sample operating-condition inputs
used by the physics-informed method (\S\ref{sec:stress}).

\paragraph{Windows, splits, and leakage control.}
Records are split at the \emph{record} level ($70/15/15$), stratified by
class$\times$severity, into $2{,}464 / 528 / 528$ records; training and
evaluation use 1\,s windows ($20{,}480$ samples, five non-overlapping windows
per record) cut at load time by a group-aware \texttt{WindowedView} that never
crosses a record boundary, giving $12{,}320 / 2{,}640 / 2{,}640$ windows. Because
all five windows of a record share its fault, severity, operating point, and
noise realization, the windows are \emph{not} independent --- a fact that
governs the statistical unit throughout (\S\ref{sec:benchmark}). Splits are
leakage-checked by record hash at generation, and both independent audits
re-verified zero cross-split overlap by content-hashing the released file. For
noise-robustness evaluation the \emph{test} split is additionally exported at
three SNR levels ($20/10/5$\,dB AWGN, calibrated to each record's clean RMS), so
frozen checkpoints are evaluated across noise without retraining.

\paragraph{Release and datasheet.}
We release the generator, the dataset
(\texttt{data/generated/dataset\_v2.h5}; SHA-256 pinned in
\texttt{results/PROVENANCE\_MANIFEST.md}), and a datasheet following
\citet{gebru2021datasheets} (Table~\ref{tab:datasheet}), under CC-BY-4.0.
Generation is deterministic (global seed 42; per-record seed $42+i$).
\S\ref{sec:repro} details the full reproducibility chain.

% T4 -- datasheet summary. Sources: dataset_card.yaml; experiments/DATASET_V2.md;
% results/PROVENANCE_MANIFEST.md (content hashes).
\begin{table}[t]
  \centering
  \caption{JBFD-11 datasheet summary (measured from the released file, not
    aspirational). Full datasheet follows \citet{gebru2021datasheets}.}
  \label{tab:datasheet}
  \begin{tabular}{ll}
    \toprule
    Property & Value \\
    \midrule
    Classes & 11 journal-bearing faults (8 single + 3 mixed) \\
    Records & 3{,}520 (320 / class; 80 / class / severity) \\
    Severity levels & 4 (incipient, mild, moderate, severe) \\
    Record length / rate & 5\,s @ 20{,}480\,Hz ($102{,}400$ samples); $\Omega_0=60$\,Hz \\
    Windows & 1\,s ($20{,}480$ samples), 5 non-overlapping per record \\
    Splits (records) & 2{,}464 / 528 / 528 (train / val / test), record-level \\
    Splits (windows) & 12{,}320 / 2{,}640 / 2{,}640 \\
    Split design & group-aware, stratified by class$\times$severity \\
    Leakage & 0 cross-split record-hash overlap (audit-verified) \\
    Operating conditions & speed $\pm10\%$, load 30--100\%, temp.\ 40--80$^\circ$C \\
    Noise / SNR variants & 7-layer noise; test at clean, 20, 10, 5\,dB \\
    Generator CI & 34 spectral-signature tests \\
    License / seed & CC-BY-4.0; global seed 42 \\
    Dataset SHA-256 & \texttt{f72ad35b\ldots afcee0} (full hash in manifest) \\
    \bottomrule
  \end{tabular}
\end{table}


% =====================================================================
\section{Benchmark \& protocol [C2]}
\label{sec:benchmark}

\paragraph{Protocol (frozen in advance).}
The benchmark task is 11-class fault classification on 1\,s windows. The
protocol was frozen before running (\texttt{experiments/PROTOCOL.md}): every deep
model is trained with an identical budget (Adam, $\mathrm{lr}=10^{-3}$; batch
64; $\le$60 epochs; early stopping with patience 10 on validation accuracy; no
learning-rate schedule; cross-entropy loss; no augmentation); model selection
uses the validation split only; and the clean \texttt{test} split and its
\texttt{test\_snr*} variants are each touched exactly once, after training. We
evaluate eleven models: eight deep networks (\texttt{cnn1d},
\texttt{attention\_cnn}, \texttt{cnn\_lstm}, \texttt{resnet18},
\texttt{patchtst}, \texttt{hybrid\_pinn}, \texttt{physics\_constrained\_cnn},
\texttt{multitask\_pinn}), three classical baselines (random forest, SVM,
gradient boosting on a 36-feature extraction), and a soft-voting ensemble; each
at three seeds.

\paragraph{The record is the inferential unit.}
The $2{,}640$ test windows derive from 528 records, five per record, and share
within-record structure; treating them as independent overstates significance
(an early error, caught by audit and corrected). We therefore aggregate each
record by soft vote (mean class score over its five windows) and compute all
statistics on the \textbf{528 records} as the independent unit. We report this
both because it is the correct estimand here and because the same
window-versus-record distinction recurs, at the seed level, in the stress-test
(\S\ref{sec:stress}).

\paragraph{Results: near-ceiling, no accuracy advantage.}
At the record level the task is near-ceiling (Table~\ref{tab:benchmark}): the
strongest classical baseline is random forest ($98.74\%$) and the strongest deep
model is \texttt{cnn\_lstm} ($99.43\%$). No model --- physics-labelled or not ---
shows an accuracy advantage. The best vanilla deep model (\texttt{cnn\_lstm},
best-validation seed $99.62\%$) and the best physics-labelled model
(\texttt{physics\_constrained\_cnn}, $99.62\%$) are tied: the paired gap is
$+0.00$ points ($95\%$ CI $[+0.00,+0.00]$, cluster bootstrap by record) with an
exact McNemar $p=1$.

\paragraph{Honest relabelling.}
The physics-labelled rows do not test physics-informed \emph{training}, and are
labelled accordingly. \texttt{physics\_constrained\_cnn} here is CE-only: its
physics term is a separate loss that the benchmark training loop never invokes,
so the row is architecture-only with the physics loss OFF.
\texttt{hybrid\_pinn} carries a rolling-element (not journal-bearing) physics
branch plus constant metadata, and \texttt{multitask\_pinn} runs single-task
with its auxiliary heads unused. The benchmark is thus a clean
\emph{classification} benchmark on which physics architecture confers no
accuracy edge; the actual test of physics-informed \emph{learning} is the
stress-test of \S\ref{sec:stress}.

% T1 -- record-level benchmark (11 models). Source:
% results/benchmark/summary_record_level.md  (every number traced there).
\begin{table}[t]
  \centering
  \caption{Record-level benchmark (528 records, soft-vote over the 5 windows
    per record; mean$\pm$std over 3 seeds). The task is near-ceiling and no
    model --- physics-labelled or not --- shows an accuracy advantage: the best
    vanilla deep model (\texttt{cnn\_lstm}) ties the CE-only
    \texttt{physics\_constrained\_cnn} at the best-validation seed (gap $+0.00$,
    exact McNemar $p=1$). Physics-labelled rows are honestly relabelled (see
    text): \texttt{physics\_constrained\_cnn} $=$ CE-only / architecture
    (physics loss OFF), \texttt{hybrid\_pinn} $=$ rolling-element branch $+$
    constant metadata, \texttt{multitask\_pinn} $=$ single-task.}
  \label{tab:benchmark}
  \begin{tabular}{llc}
    \toprule
    Model & Record acc.\ (\%) & Note \\
    \midrule
    random\_forest          & $98.74 \pm 0.09$ & classical \\
    svm                     & $97.73 \pm 0.00$ & classical \\
    gradient\_boosting      & $97.92 \pm 0.15$ & classical \\
    \midrule
    cnn1d                   & $93.43 \pm 4.26$ & \\
    attention\_cnn          & $93.94 \pm 4.27$ & \\
    cnn\_lstm               & $\mathbf{99.43 \pm 0.15}$ & best vanilla deep \\
    resnet18                & $99.18 \pm 0.09$ & \\
    patchtst                & $90.40 \pm 0.47$ & \\
    \midrule
    physics\_constrained\_cnn & $98.99 \pm 0.45$ & physics-labelled: CE-only (loss OFF) \\
    hybrid\_pinn            & $90.34 \pm 0.94$ & physics-labelled: rolling-element $+$ const.\ meta \\
    multitask\_pinn         & $90.47 \pm 0.62$ & physics-labelled: single-task \\
    \bottomrule
  \end{tabular}
\end{table}


% =====================================================================
\section{Stress-testing the tested physics-informed methods [C3]}
\label{sec:stress}
\textit{[TODO C3] The implemented method: PC-CNN + band-energy-vs-healthy-
reference consistency loss (per-sample rpm). The five regimes. The
pre-registered \S8.8 $n{=}12$ grid (CE-only / correct / scramble / random-band
$\times$ 12 seeds; Table~\ref{tab:grid}, Fig.~\ref{fig:spread}); the seed-level
estimand + fixed decision rule; the result (no arm clears the bar). The other
regimes (\S8.2 data-efficiency, \S8.3 severity-OOD, \S8.6a interpretability,
\S8.6b calibration), each negative (Table~\ref{tab:regimes}).}

% T2 -- the pre-registered S8.8 n=12 grid (the headline). Source:
% results/noise_seed_robustness/noise_seed_robustness_record_level.json  (every number traced).
\begin{table}[t]
  \centering
  \caption{Pre-registered grid (protocol~\S8.8), $n{=}12$: per-arm clean$\rightarrow$5\,dB
    record-level degradation (mean$\pm$std and median, pts), robust-seed count
    (degradation $<\tau=1.0$\,pt), and the \emph{seed-level} Wilcoxon
    signed-rank $p$ versus CE-only. No arm clears the pre-registered bar
    ($\ge$10/12 robust \emph{and} $p<0.05$): correct physics is
    indistinguishable from CE-only, and the non-physical random-band control is
    numerically the most robust yet still fails. The $n{=}3$ positive does not
    replicate. The representative-seed McNemar is deliberately \emph{not}
    reported here as evidence (see \S\ref{sec:caution}).}
  \label{tab:grid}
  \begin{tabular}{lccccc}
    \toprule
    Arm & $w$ & degr (mean$\pm$std) & median & robust & Wilcoxon $p$ \\
    \midrule
    CE-only        & 0   & $3.54 \pm 3.85$ & 2.46 & 4/12 & --- (ref.) \\
    correct        & 1.0 & $3.47 \pm 3.42$ & 2.84 & 5/12 & $0.79$ \\
    scramble       & 1.0 & $4.89 \pm 6.08$ & 0.47 & 7/12 & $0.75$ \\
    random-band    & 1.0 & $2.15 \pm 3.34$ & 0.38 & 9/12 & $0.21$ \\
    \bottomrule
  \end{tabular}
\end{table}

% T3 -- the other regimes at a glance (each negative, record-level). Sources:
% results/phase5_bandenergy/summary_record_level.json (8.2 data-eff, 8.3 OOD);
% results/xai_alignment/alignment.json (8.6a); results/uncertainty/calibration.json (8.6b)
% [STUB: populated when Section 5 (Stress-testing, C3) is drafted.]
\begin{table}[t]
  \centering
  \caption{\textit{[TODO T3]} The other stress-test regimes, each record-level
    and each negative: data-efficiency (\S8.2), severity-OOD (\S8.3),
    interpretability (\S8.6a), calibration (\S8.6b).}
  \label{tab:regimes}
  \begin{tabular}{lll}
    \toprule
    Regime & Result & Verdict \\
    \midrule
    \textit{[TODO: 4 rows]} & --- & negative \\
    \bottomrule
  \end{tabular}
\end{table}


\begin{figure}[t]
  \centering
  \figfromfile{F1_degradation_spread.pdf}
  \caption{\textit{[TODO caption F1]} Per-seed clean$\rightarrow$5\,dB
    record-level degradation (528 records, soft-vote) for the four \S8.8 arms,
    $n{=}12$ seeds; the dashed line is the pre-registered robustness threshold
    $\tau=1.0$\,pt. The spread that fooled $n{=}3$ is visible. Source:
    \texttt{results/p7\_strengthen/p7\_strengthen\_record\_level.json}.}
  \label{fig:spread}
\end{figure}

% =====================================================================
\section{Methodological caution [C4]}
\label{sec:caution}
\textit{[TODO C4] The $n{=}3\!\rightarrow\!n{=}12$ dissolution
(Fig.~\ref{fig:dissolution}): a within-seed McNemar of $14$--$0$
($p=1.2\times10^{-4}$) at the representative seed becomes a seed-level Wilcoxon
$p=0.79$ at $n{=}12$; the random non-fault-band control is the most robust arm.
The broken-loss saga (five defects, all passing a green suite).
Recommendations: seed as the inferential unit; pre-register the estimand;
scrambled + matched-strength non-physics controls.}

\begin{figure}[t]
  \centering
  \figfromfile{F2_n3_to_n12_dissolution.pdf}
  \caption{\textit{[TODO caption F2 --- HEADLINE]} The apparent
    noise-robustness benefit is a seed artifact. \emph{Left}: at $n{=}3$ (seeds
    0--2) the correct-physics arm appears far more robust than CE-only.
    \emph{Right}: the same comparison under the pre-registered $n{=}12$
    seed-level test is a null (Wilcoxon $p=0.79$). Same data, same code path;
    the gap dissolves as seeds grow. Source:
    \texttt{results/p7\_strengthen/p7\_strengthen\_record\_level.json}.}
  \label{fig:dissolution}
\end{figure}

% =====================================================================
\section{Reproducibility [C5]}
\label{sec:repro}
\textit{[TODO C5] results/PROVENANCE\_MANIFEST.md; content hashes; the Zenodo
checkpoint archive (DOI pending --- ``will archive''); the independent
from-checkpoint reproduction (audit\_reports/$\ldots$2026-06-24$\ldots$);
hosting / license (CC-BY-4.0) / maintenance plan; D\&B reproducibility
checklist.}

% =====================================================================
\section{Limitations \& scope}
\label{sec:limitations}
\textit{[TODO] The section reviewers will read: synthetic-only / no real rig;
near-ceiling clean task; single dataset; the negative is about the \emph{methods
we implemented and tested}, not a claim that no physics prior could ever help.}

% =====================================================================
\section{Conclusion}
\label{sec:conclusion}
\textit{[TODO]}

% =====================================================================
\bibliographystyle{plainnat}
\bibliography{references}

% =====================================================================
\appendix
\section{Deployment}
\label{app:deployment}
\textit{[TODO C5-appendix] ResNet18 $\rightarrow$ ONNX FP32 $\approx$13\,ms/window
CPU; INT8 4$\times$ smaller but 10--15$\times$ slower (honest negative). Source:
\texttt{results/deployment/appendix.md}.}

\end{document}
